Medical foundation models pre-trained on images and image--text pairs are reused as frozen feature
extractors, zero-shot classifiers and report generators, so their internal representations matter. This survey reviews work that analyses, decomposes, traces and
intervenes on them. It organises methods by
family, locus and concept provenance and codes \Nclaims{} records from \Ncore{} core studies into \Fnrows{}
evidence statements, each read against \Vnclasses{} classes of validity test that state what a test
licenses \emph{and what it does not}. Counts describe this corpus, \Pctpreprint\% of it preprints, not the literature. Clinical findings, anatomy, demographic and acquisition
attributes are reported as recoverable from frozen encoders in four to seven modalities, but controls on
probe validity are rare, few targeted interventions report a matched control and no study meets the criterion
stated here for a concept-specific causal contribution, so evidence that models use what is recoverable
remains limited. Centre, scanner and acquisition
structure is reported recoverable in every record testing it but one, and whether removing it restores
generalisation is operation-specific. Two thirds of the dictionary studies report a semantic check of their
features, none cross-seed identity. In the same-data
comparisons reviewed, selected geometric and shortcut properties are reported in both medical and general
models, and because no comparison holds pre-training data, objective, architecture and scale fixed, the effect of
medical pre-training is unidentified.
The connector output and text encoder are almost unmeasured in this corpus. The coded corpus is released.
